\documentclass[11pt]{article}
\usepackage[margin=1in]{geometry}
\usepackage{amsmath,amssymb,mathtools,amsthm}
\usepackage{booktabs,array,longtable}
\usepackage{enumitem}
\usepackage{hyperref}
\usepackage{xcolor}

\newcommand{\SB}{\textsc{ScaleBridge}}
\newcommand{\R}{\mathbb{R}}
\newcommand{\Hspace}{\mathcal{H}}
\newcommand{\Study}{\mathcal{S}}
\newcommand{\Trial}{\tau}
\newcommand{\Data}{\mathcal{D}}
\newcommand{\Provider}{\mathfrak{P}}
\newcommand{\Fingerprint}{\mathfrak{F}}
\newcommand{\Freeze}{\mathcal{A}^{\star}}
\newtheorem{invariant}{Invariant}
\newtheorem{proposition}{Proposition}

\title{\SB{} Phase E0 --- E0-8\\
Generic Hyperparameter Optimization, Study Tracking,\\
and Frozen-Configuration Framework\\
Mathematical and Scientific Contract v1}
\author{ScaleBridge Research Framework}
\date{Ratified August 2026}

\begin{document}
\setlength{\emergencystretch}{3em}
\maketitle

\begin{abstract}
This document freezes the E0-8 mathematical, scientific, and reproducibility
contract for generic hyperparameter optimization (HPO) in \SB{}. E0-8 is
method-neutral. It owns the lifecycle of a study, trial orchestration, Optuna
integration, MLflow tracking, TRAIN-only data-leakage enforcement, study
fingerprints, compatible resume, failure/pruning semantics, single- and
multi-objective bookkeeping, and production of a frozen hyperparameter
configuration. E0-8 does \emph{not} define method-specific hyperparameters,
search-space meanings, model-fitting equations, objective science,
representative-subset science, pruning validity, or final-model fitting. Those
responsibilities belong to the future E.1/E.2/E.3/E.4 method families.

Let $m$ denote a method and $h\in\Hspace_m$ its hyperparameter vector. A method
provider supplies the search space, objective contract, TRAIN-side data policy,
optional pruning interface, resource/repetition metadata, and deterministic
final-selection policy. E0-8 then orchestrates trials of the form
\[
M_j=\operatorname{Fit}_m(\Data_{\mathrm{fit},j};h_j,s_j),
\qquad
\mathbf y_j=\mathbf J_m(M_j,\Data_{\mathrm{score},j}),
\]
while keeping $h$ distinct from fitted model parameters such as $\Theta$,
$\psi$, or a posterior representation $\Pi_\Theta$. Hyperparameter selection
is restricted to Phase-D TRAIN data. Phase-D VALIDATION and TEST are excluded
from all HPO selection influence; any fit/score split or cross-validation used
inside a trial must be derived from TRAIN only.

Optuna is the study/trial optimization authority and MLflow is the
experiment/provenance authority. This division matches Optuna's study/trial,
sampler, pruner, persistence, and dynamic-search-space abstractions
\cite{AkibaEtAl2019Optuna,OptunaDocs2026} and MLflow's experiment/run tracking
of parameters, metrics, metadata, and artifacts \cite{MLflowDocs2026}. The
precise ownership boundary, however, is a \SB{}-specific contract. General HPO
notation and reproducibility considerations are consistent with the broader HPO
literature \cite{BischlEtAl2023HPO}.
\end{abstract}

% =============================================================================
\section{Scope and responsibility boundary}

E0-8 sits between method-specific HPO declarations and method-specific final
model fitting:
\[
\boxed{
\text{E.x HPO provider}
\rightarrow
\text{E0-8 study}
\rightarrow
h^{\star}_{\mathrm{frozen}}
\rightarrow
\text{E.x final fit}
\rightarrow
\text{E0-7 bundle}.
}
\]
Here E.x denotes one of E.1/E.2/E.3/E.4.

\begin{invariant}[Generic orchestration]
E0-8 owns generic HPO orchestration and tracking. E.1/E.2/E.3/E.4 own the
scientific meaning and declaration of all method-specific hyperparameters and
objectives.
\end{invariant}

Accordingly, E0-8 must not contain method branches such as ``if Base-PINODE then
use these hyperparameters'' or ``if Bayesian then optimize these sampler
settings.'' Instead, each method registers a provider that satisfies the common
E0-8 interface.

\paragraph{Explicit non-responsibilities.}
E0-8 does not own model equations, forward discretization, backpropagation,
CasADi/IPOPT parameter estimation, Bayesian posterior inference, neural-network
architecture science, scientific loss design, final production-model fitting,
Sim1/Sim2/Sim3 evaluation, MPC, RL, or SmartBuildingsSim. E0-10 remains the
authoritative comprehensive real testing campaign after the method families
exist.

% =============================================================================
\section{Hyperparameters versus fitted model parameters}

For method $m$, let
\[
h\in\Hspace_m
\]
be the hyperparameter vector. Let the fitted model quantities be represented
generically by
\[
\vartheta_m\in\mathcal V_m,
\]
where $\vartheta_m$ may contain physical parameters $\Theta$, neural weights
$\psi$, a posterior object $\Pi_\Theta$, noise parameters, or other
method-owned fitted quantities.

\begin{invariant}[Coordinate separation]
\[
\boxed{h\neq\vartheta_m.}
\]
E0-8 chooses and freezes $h$; the method-specific fitter estimates
$\vartheta_m$ conditional on $h$.
\end{invariant}

For trial $j$,
\[
\vartheta_{m,j}^{\star}
=
\operatorname{Fit}_m(\Data_{\mathrm{fit},j};h_j,s_j),
\]
and the trial score is
\[
\mathbf y_j
=
\mathbf J_m(\vartheta_{m,j}^{\star},\Data_{\mathrm{score},j}).
\]
E0-8 invokes these method-provided operations but does not redefine their
internal mathematics.

% =============================================================================
\section{Method HPO provider contract}

Each method $m$ provides a generic declaration
\[
\boxed{
\Provider_m=
(\mathcal S_m,\mathcal J_m,\mathcal D_m,\mathcal Q_m,\mathcal R_m,\mathcal F_m),
}
\]
with:
\begin{align*}
\mathcal S_m &: \text{search-space declaration},\\
\mathcal J_m &: \text{objective names, directions, and evaluator},\\
\mathcal D_m &: \text{TRAIN-side HPO data-selection policy},\\
\mathcal Q_m &: \text{optional pruning/intermediate-reporting contract},\\
\mathcal R_m &: \text{resource, repetition, timeout, and seed requirements},\\
\mathcal F_m &: \text{deterministic final hyperparameter-selection policy}.
\end{align*}

E0-8 understands only generic domain types such as continuous/log-continuous,
integer, categorical, Boolean, and conditional parameters. The names, bounds,
categories, conditions, defaults, and scientific meanings are method-owned.
Optuna's define-by-run design motivates support for dynamic/conditional spaces
\cite{AkibaEtAl2019Optuna}, but the exact ScaleBridge provider boundary above is
ScaleBridge-specific.

\begin{invariant}[Search-space ownership]
A change to a method's hyperparameters changes the E.x provider declaration,
not the common E0-8 engine.
\end{invariant}

% =============================================================================
\section{Study object and immutable identity}

A canonical study is
\[
\boxed{
\Study=
(I,m,V_m,\mathcal S_m,\mathcal J_m,\mathcal D_m,s_{\mathrm{study}},
\mathcal O,\mathcal T),
}
\]
where $I$ is study identity, $m$ method identity, $V_m$ provider/search-space
version, $s_{\mathrm{study}}$ study seed, $\mathcal O$ generic Optuna
configuration, and $\mathcal T$ tracking configuration.

The study stores a frozen snapshot of the provider's search-space and objective
contract at creation time. Subsequent source-code changes to a provider do not
silently reinterpret already recorded trials.

\paragraph{Study fingerprint.}
Define the immutable resume fingerprint
\[
\boxed{
\Fingerprint(\Study)
=
H\!\left(
I,m,V_m,\operatorname{canon}(\mathcal S_m),
\operatorname{canon}(\mathcal J_m),
\operatorname{fp}(\mathcal D_m),
s_{\mathrm{study}},v_{\mathrm{E0-8}}
\right),
}
\]
where $H$ is a stable cryptographic digest and $\operatorname{canon}$ denotes a
deterministic canonical serialization.

\begin{invariant}[Compatible resume only]
A stored study may be resumed only when
\[
\Fingerprint_{\mathrm{stored}}=\Fingerprint_{\mathrm{requested}}.
\]
If the method version, search-space snapshot, objective contract, HPO data
selection, seed policy, or E0-8 schema is incompatible, resume fails clearly
rather than appending scientifically incomparable trials.
\end{invariant}

% =============================================================================
\section{Phase-D TRAIN-only HPO data authority}

Let the authoritative Phase-D outer partition be
\[
\Data
=
\Data_{\mathrm{train}}
\cup
\Data_{\mathrm{validation}}
\cup
\Data_{\mathrm{test}}
\cup
\Data_{\mathrm{excluded}},
\]
with disjoint primary partitions.

The HPO source satisfies
\[
\boxed{
\Data_{\mathrm{HPO}}\subseteq\Data_{\mathrm{train}}.
}
\]
Therefore
\[
\Data_{\mathrm{HPO}}\cap\Data_{\mathrm{validation}}=\varnothing,
\qquad
\Data_{\mathrm{HPO}}\cap\Data_{\mathrm{test}}=\varnothing.
\]

\begin{invariant}[No validation/test leakage into HPO]
Phase-D VALIDATION and TEST cannot affect hyperparameter proposals, pruning,
trial scores, ranking, Pareto construction, or final hyperparameter selection.
\end{invariant}

If a method requires fit/score separation during HPO, it must derive it from
TRAIN only. For example,
\[
\Data_{\mathrm{train}}
=
\Data_{\mathrm{inner\mbox{-}fit}}
\dot\cup
\Data_{\mathrm{inner\mbox{-}score}}
\dot\cup
\Data_{\mathrm{unused}},
\]
or a deterministic TRAIN-only cross-validation construction may be used. E0-8
validates the outer-partition boundary and records the method-provided inner
selection.

% =============================================================================
\section{Representative-subset HPO}

Expensive methods may declare
\[
\Data_{\mathrm{rep}}\subseteq\Data_{\mathrm{train}}.
\]
The science of representative selection belongs to E.x. E0-8 records, at
minimum, the policy identifier/version, selected cases/zones/time ranges,
selection seed when relevant, and a deterministic data-selection fingerprint.
The trial objective is then evaluated on the provider-declared TRAIN-side HPO
construction, for example
\[
\mathbf J_m(h;\Data_{\mathrm{rep}}).
\]

\begin{invariant}[Selection provenance]
E0-8 never silently chooses which buildings, zones, seasons, or time blocks are
``representative.'' It persists exactly what the method provider selected and
why that selection is reproducible.
\end{invariant}

% =============================================================================
\section{Trial contract}

Trial $j$ is
\[
\boxed{
\Trial_j=
(j,h_j,s_j,\mathbf y_j,q_j,a_j,e_j,t_j),
}
\]
where $q_j$ is trial state, $a_j$ artifact/metadata references, $e_j$ failure or
pruning evidence when applicable, and $t_j$ timing/resource metadata.

A deterministic trial seed is derived from the study seed and trial identity:
\[
\boxed{
s_j=\operatorname{SeedDerive}(s_{\mathrm{study}},j,I).
}
\]
The concrete derivation must be stable across supported machines and Python
processes; implementation must not depend on randomized language-runtime hash
seeds.

The method receives $h_j$, $s_j$, and the authorized TRAIN-side data selection,
then returns objective value(s), intermediate reports if supported, artifacts,
and terminal state.

% =============================================================================
\section{Trial terminal states and failure semantics}

E0-8 distinguishes at least
\[
q_j\in\{\mathrm{COMPLETE},\mathrm{PRUNED},\mathrm{FAILED}\}.
\]

\begin{invariant}[State distinction]
\[
\boxed{
\mathrm{COMPLETE}\neq\mathrm{PRUNED}\neq\mathrm{FAILED}.
}
\]
A failed trial is not converted into an artificially poor objective value, and
a pruned trial is not reported as a completed scientific evaluation.
\end{invariant}

Trial-level failures retain the sampled hyperparameters, trial seed, exception
class/message, traceback or structured failure evidence, machine/environment
identity, and elapsed/runtime metadata. A configured study may continue after a
trial-level failure. Study/storage corruption, incompatible resume, leakage
violation, or invalid provider contract is a study-level failure.

% =============================================================================
\section{Pruning contract}

If method $m$ declares pruning scientifically meaningful, it may expose
intermediate values
\[
\mathbf y_j^{(1)},\mathbf y_j^{(2)},\ldots,\mathbf y_j^{(r)}.
\]
E0-8 reports them to the configured Optuna pruning mechanism and records the
pruning decision and step. Optuna explicitly separates study, sampler, and
pruner concepts \cite{OptunaDocs2026}.

\begin{invariant}[Method-declared pruning compatibility]
E0-8 never assumes that epochs, optimizer iterations, IPOPT iterations, MCMC
iterations, or Bayesian filtering steps are comparable intermediate HPO scores.
The method provider must explicitly authorize and define intermediate reporting.
\end{invariant}

% =============================================================================
\section{Single-objective selection}

For a scalar objective $J$ with direction $d_m$, the completed-trial set is
\[
\mathcal C=\{j:q_j=\mathrm{COMPLETE}\}.
\]
For minimization,
\[
\boxed{
h^{\star}=h_{j^{\star}},\qquad
j^{\star}=\arg\min_{j\in\mathcal C}J(h_j),
}
\]
and for maximization,
\[
\boxed{
j^{\star}=\arg\max_{j\in\mathcal C}J(h_j).}
\]
The objective direction is method-declared. E0-8 does not assume that lower is
universally better.

If repeated stochastic evaluations are required for one $h_j$, the method
provider declares the repetitions and the deterministic reduction rule that
produces the trial objective. E0-8 tracks each repetition's seed/evidence but
does not invent the scientific reducer.

% =============================================================================
\section{Multi-objective selection}

For $q>1$ objectives,
\[
\mathbf J(h)=[J_1(h),\ldots,J_q(h)]^T.
\]
E0-8 and Optuna may maintain the nondominated/Pareto set
\[
\boxed{
\mathcal P^{\star}
=\operatorname{Pareto}\{(h_j,\mathbf J(h_j)):j\in\mathcal C\}.
}
\]

\begin{invariant}[No invented scalarization]
E0-8 may compute and persist the Pareto set but does not invent a scientific
scalarization or utility function. If one final $h^{\star}$ is required, E.x
must provide a deterministic final-selection rule $\mathcal F_m$. If no such
rule exists, E0-8 returns the Pareto set without falsely declaring one trial
``best.''
\end{invariant}

% =============================================================================
\section{Optuna orchestration authority}

Optuna provides the optimization-layer concepts used by E0-8: studies and
trials, persistent storage, samplers, pruners, single/multiple directions, and
dynamic search spaces \cite{AkibaEtAl2019Optuna,OptunaDocs2026}. E0-8 wraps
these capabilities in ScaleBridge-specific scientific checks.

E0-8 owns generic registries/configuration for compatible sampler, pruner, and
storage capabilities. E.x may restrict compatibility; for example, a method may
declare pruning disabled or require a particular objective structure. E0-8
does not redefine Optuna's optimization algorithm internally.

\begin{invariant}[Optimization authority]
\[
\boxed{\text{Optuna is the HPO optimization authority.}}
\]
The E0-8 layer adds scientific identity, leakage protection, reproducibility,
provenance, and output contracts around that authority.
\end{invariant}

% =============================================================================
\section{MLflow tracking authority}

MLflow Tracking organizes executions as runs grouped by experiments and records
parameters, metrics, metadata, and artifacts \cite{MLflowDocs2026}. E0-8 uses a
study-parent / trial-child logical hierarchy:
\[
\boxed{
\text{HPO study parent run}
\rightarrow
\{\text{trial runs}\}
\rightarrow
\text{study finalization evidence}.
}
\]

The study parent records method/provider identity, Phase-D lineage/data
fingerprint, search-space snapshot, objective contract, study seed, Optuna
sampler/pruner/storage identity, environment, and study fingerprint. Each trial
records sampled hyperparameters, objective values, trial seed/state, duration,
method metrics, selected artifacts, and failure/pruning evidence. Finalization
records the frozen configuration or Pareto/selection decision.

\begin{invariant}[Tracking authority]
\[
\boxed{\text{MLflow is the experiment/provenance authority.}}
\]
MLflow does not replace Optuna's study/trial optimization logic, and E0-8 does
not use MLflow run ordering as an optimization algorithm.
\end{invariant}

% =============================================================================
\section{Frozen hyperparameter configuration}

The principal downstream artifact is
\[
\boxed{
\Freeze=
(h^{\star},I,j^{\star},\Fingerprint,\operatorname{fp}(\Data_{\mathrm{HPO}}),
\operatorname{fp}(\mathcal S_m),\operatorname{fp}(\mathcal J_m),\rho),
}
\]
where $\rho$ contains provenance such as MLflow run identity, Optuna study
identity, method/provider version, and selection evidence.

For a Pareto study without one final selection, the corresponding frozen output
contains the Pareto trial set and its complete provenance rather than a
fabricated scalar optimum.

\begin{invariant}[Freeze, do not final-fit]
E0-8 freezes $h^{\star}$ (or the Pareto result) but does not fit the final
production model. E.x consumes the frozen configuration, performs the final
allowed fit, and subsequently exports the fitted model to E0-7.
\end{invariant}

% =============================================================================
\section{Standard E0-8 artifact contract}

The generic study layer should produce at least the following logical artifacts:
\begin{center}
\begin{tabular}{ll}
\toprule
Artifact & Purpose\\
\midrule
\texttt{study\_manifest.json} & immutable study identity and schema\\
\texttt{search\_space\_snapshot.json} & exact provider search-space snapshot\\
\texttt{objective\_contract.json} & names, directions, evaluator identity\\
\texttt{data\_selection\_manifest.json} & TRAIN-only HPO selection + fingerprint\\
\texttt{trials.parquet} & trial parameters, objectives, state, seed, timing\\
\texttt{study\_summary.json} & aggregate study result/status\\
\texttt{frozen\_hyperparameters.json} & selected $h^{\star}$ + provenance\\
\bottomrule
\end{tabular}
\end{center}
For multi-objective studies, \texttt{pareto\_trials.parquet} and a
\texttt{selection\_manifest.json} may replace or complement the scalar
best-trial artifact. Exact filesystem placement is an implementation concern;
the semantic contents above are the authority.

% =============================================================================
\section{Interaction with E0-7, E0-9, and E0-10}

The eventual E0-7 portable model bundle should preserve the HPO lineage of the
finished model, including at least the E0-8 study identity, selected trial or
selection manifest, frozen hyperparameter payload or its content hash, and
method/provider version. E0-7 never reruns HPO.

E0-9 validates E0-8 infrastructure/readiness across qualified machines and
repairs environment discrepancies without changing the HPO science. E0-10 is
the authoritative comprehensive real testing campaign and, once E.1--E.4
methods exist, exercises the chain
\[
\boxed{
\text{Phase D}
\rightarrow
\text{E.x provider}
\rightarrow
\text{E0-8 study}
\rightarrow
h^{\star}
\rightarrow
\text{E.x final fit}
\rightarrow
\text{E0-7 bundle}
\rightarrow
\text{evaluation}.
}
\]

% =============================================================================
\section{Worked generic example}

This example is intentionally synthetic and does not freeze any real E.1--E.4
hyperparameter names. Suppose a future method $m_{\mathrm{demo}}$ declares
three hyperparameters:
\[
h=(h_c,h_i,h_k),
\]
where $h_c$ is continuous, $h_i$ integer, and $h_k$ categorical. The provider
might declare
\[
h_c\in[10^{-4},10^{-2}]\ \text{on a log scale},\qquad
h_i\in\{2,3,4\},\qquad
h_k\in\{A,B\}.
\]
These names and ranges are illustrative only; a real E.x provider owns its own
space.

Assume Phase D contains 100 TRAIN windows, 20 VALIDATION windows, and 20 TEST
windows. The method declares an HPO representative subset of 40 TRAIN windows
and a deterministic inner split
\[
30\ \text{TRAIN windows for trial fitting},\qquad
10\ \text{TRAIN windows for trial scoring}.
\]
No VALIDATION or TEST window is eligible for the study.

With study seed $s_{\mathrm{study}}=1234$, trial 7 receives
\[
s_7=\operatorname{SeedDerive}(1234,7,I).
\]
Suppose Optuna proposes
\[
h_7=(3\times10^{-3},3,B).
\]
The method performs
\[
\vartheta_7^{\star}
=\operatorname{Fit}_{m_{\mathrm{demo}}}(\Data_{30};h_7,s_7)
\]
and returns a minimization score
\[
J(h_7)=0.184
\]
on the authorized 10-window inner score set. E0-8 records $h_7$, $s_7$,
$J(h_7)$, state COMPLETE, timing/resources, MLflow trial identity, and the data
fingerprint. If trial 8 raises a numerical exception, it remains FAILED rather
than receiving an invented score. If trial 9 is pruned under a provider-enabled
pruning contract, it remains PRUNED.

After all trials, suppose completed objectives are
\[
0.231,\quad 0.184,\quad 0.201.
\]
For a minimization study,
\[
h^{\star}=h_7.
\]
E0-8 writes \texttt{frozen\_hyperparameters.json} containing $h_7$ and full
study/data/search-space/objective provenance. E0-8 does not train the final
production model. The method later consumes $h_7$, performs its final fit under
its own E.x contract, and sends the fitted model to E0-7.

% =============================================================================
\section{Qualification requirements}

E0-8 implementation qualification must include, at minimum:
\begin{enumerate}[label=\arabic*.]
\item provider-agnostic search-space registration with no method-specific branch;
\item TRAIN-only enforcement rejecting Phase-D VALIDATION/TEST leakage;
\item deterministic data-selection fingerprints and trial seeds;
\item single-objective minimization and maximization studies;
\item multi-objective Pareto persistence and method-owned final selection;
\item COMPLETE/PRUNED/FAILED distinction;
\item compatible resume with exact fingerprint match and incompatible-resume rejection;
\item persistent Optuna study recovery;
\item MLflow study/trial/finalization tracking linkage;
\item deterministic frozen-hyperparameter artifact creation;
\item failure recovery without corrupting completed trial history;
\item downstream handoff semantics to E.x and HPO-lineage preservation for E0-7.
\end{enumerate}
Actual method-family search spaces are not required for E0-8 framework
qualification because E.1--E.4 are not yet implemented; synthetic providers are
the correct E0-8 framework fixture until those method families exist.

% =============================================================================
\section{Ratified E0-8 contract}

The following decisions are frozen:
\begin{enumerate}[label=\arabic*.]
\item E0-8 is generic and method-neutral.
\item E.1/E.2/E.3/E.4 define hyperparameter names, domains, conditional dependencies, and meanings.
\item E.1--E.4 define trial objective evaluators and objective directions.
\item E.1--E.4 define whether pruning is scientifically meaningful.
\item E.1--E.4 define any representative HPO subset science.
\item E0-8 never hard-codes method-specific hyperparameters.
\item Hyperparameters $h$ remain distinct from fitted quantities $\Theta$, $\psi$, and $\Pi_\Theta$.
\item E0-8 orchestrates trials; the method performs trial fitting.
\item HPO consumes only Phase-D TRAIN data.
\item Phase-D VALIDATION and TEST cannot influence hyperparameter selection.
\item Any inner fit/score split or cross-validation is derived from TRAIN only.
\item Exact HPO data selection and fingerprint are persisted.
\item Single- and multi-objective studies are supported generically.
\item E0-8 persists Pareto fronts but never invents method-specific scalarization.
\item A method requiring one final multi-objective configuration supplies deterministic selection logic.
\item COMPLETE, PRUNED, and FAILED remain distinct.
\item Study/trial seeds and reproducibility metadata are first-class.
\item Resume requires immutable study-fingerprint equality.
\item Optuna is the HPO optimization authority.
\item MLflow is the experiment/provenance authority.
\item E0-8 standardizes study/trial/frozen-configuration artifacts.
\item E0-8 freezes $h^{\star}$ but does not own final production-model fitting.
\item E.x consumes $h^{\star}$, fits the final model, and hands it to E0-7.
\item E0-10 remains the authoritative real end-to-end testing campaign.
\end{enumerate}

\bibliographystyle{plain}
\bibliography{ScaleBridge_Thermal_Modeling_References}

\end{document}
